\chapter{Voorwoord}

Dit boek heeft een eenvoudige ambitie: één samenhangend scheikundeboek dat
alles bestrijkt van het eerste schooljaar tot het einde van de bachelor,
geschreven in het Nederlands, voor lezers waar ook ter wereld.

De stijl is beknopt en streng: lesteksten opgebouwd uit definities,
voorbeelden, proposities, methoden en schema's, met zorgvuldige
redeneringen zodra die op het betreffende niveau toegankelijk zijn,
gevolgd door zorgvuldig gekozen oefeningen. Resultaten die zonder
volledige afleiding worden gegeven, zijn uitdrukkelijk als aangenomen
gemarkeerd. Van alle oefeningen staan achter in het boek volledige
uitwerkingen. Reactievergelijkingen kloppen altijd, en elk getal dat voor
een echte stof wordt genoemd, komt uit een naslagtabel.

De cursus bestaat uit vier boeken. Het eerste bestrijkt alle twaalf
schooljaren in één deel en behandelt nooit twee keer hetzelfde begrip:
elk idee wordt één keer uitgelegd, in het jaar waarin het voor het eerst
past, en een later hoofdstuk dat het nodig heeft, begint met een korte
herinnering, \emph{Wat je al weet}, voordat het verder gaat. De andere
drie boeken zijn de drie universiteitsjaren. In het eerste boek is elk
schooljaar een deel, verdeeld in hoofdstukken; elk universiteitsboek is
één jaar, verdeeld in hoofdstukken. Hoe vroeger het jaar, hoe jonger de
lezer voor wie het geschreven is: meer illustraties, meer uitgewerkte
tussenstappen en mildere oefeningen.

\section*{Hoe je dit boek leest}

De uitspraken worden per hoofdstuk genummerd en zijn met kleur gecodeerd:
definities in blauw, stellingen in rood, proposities en verwanten in
oranje, en methoden --- korte recepten voor standaardtaken --- in groen.
De oefeningen lopen van ($\star$) voor rechtstreekse toepassingen tot
($\star\star\star$) voor lastigere opgaven. Omkaderde blokken bevatten
wat niet bij de redenering zelf hoort: een herinnering aan wat een eerder
hoofdstuk heeft uitgelegd, hoe een handeling in een laboratorium wordt
uitgevoerd, een bladzijde geschiedenis, of de gevaren van een product.

\section*{Veiligheid}

Scheikunde leer je met echte stoffen, en sommige daarvan zijn gevaarlijk.
De proeven die in dit boek worden beschreven, worden uitgevoerd in een
school- of universiteitslaboratorium, door of samen met een docent, met
de beschermingsmiddelen die het laboratorium ter beschikking stelt. Geen
ervan is bedoeld om thuis te doen.

\section*{Meewerken}

De bron van dit boek staat in een openbare git-repository:
\begin{center}
\url{https://github.com/bvirrion/one-chemistry-book}
\end{center}
Bijdragen --- nieuwe hoofdstukken, correcties, betere uitleg, extra
oefeningen --- zijn welkom; zie \texttt{CONTRIBUTING.md} in de hoofdmap van
de repository.
